% ============================================================
% Section 138: n型GaAs的霍尔效应测量
% ============================================================
\section{半导体物理：n型GaAs的霍尔效应测量}
\label{sec:gaas-hall-effect}

\begin{modelbox}[题目：n型GaAs霍尔系数与迁移率]
对厚为 $0.08\,\text{cm}$ 的 n 型 GaAs，沿 $x$ 方向通以 $50\,\text{mA}$ 的电流，沿 $z$ 方向加 $0.5\,\text{T}$ 的磁场，得到 $0.4\,\text{mV}$ 的霍尔电压。求：

(1) 霍尔系数；

(2) 载流子浓度；

(3) 如材料电阻率为 $1.5\times10^{-3}\,\Omega\cdot\text{cm}$，求载流子霍尔迁移率。
\end{modelbox}

\begin{tipbox}[考点快速分析]
\begin{itemize}
    \item \textbf{霍尔电压}：$V_H=R_HIB/d$，$d$ 为沿磁场方向的样品厚度
    \item \textbf{n型霍尔系数}：$R_H=-1/(ne)<0$
    \item \textbf{载流子浓度}：$n=-1/(R_He)$
    \item \textbf{霍尔迁移率}：$\mu_H=|R_H|\sigma=|R_H|/\rho$
\end{itemize}
\end{tipbox}

\begin{derivationbox}[标准解题过程]
\textbf{(1) 霍尔系数 $R_H$}

霍尔电压与电流、磁场、样品厚度的关系：
\[
V_H=\frac{R_HIB}{d}\implies R_H=\frac{V_Hd}{IB}.
\]
已知 $V_H=0.4\,\text{mV}=4\times10^{-4}\,\text{V}$，$d=0.08\,\text{cm}=8\times10^{-4}\,\text{m}$，$I=50\,\text{mA}=0.05\,\text{A}$，$B=0.5\,\text{T}$。

n型半导体霍尔系数为负：
\begin{equation}
\boxed{R_H=-\frac{4\times10^{-4}\times8\times10^{-4}}{0.05\times0.5}\approx-1.28\times10^{-5}\,\text{m}^3/\text{C}.}
\end{equation}

\textbf{(2) 载流子浓度 $n$}

\[
n=-\frac{1}{R_He}=\frac{1}{1.28\times10^{-5}\times1.6\times10^{-19}}\approx4.88\times10^{23}\,\text{m}^{-3}.
\]
\begin{equation}
\boxed{n\approx4.88\times10^{17}\,\text{cm}^{-3}.}
\end{equation}

\textbf{(3) 霍尔迁移率 $\mu_H$}

霍尔迁移率定义为霍尔系数与电导率的乘积：
\[
\mu_H=|R_H|\sigma=\frac{|R_H|}{\rho}.
\]
已知 $\rho=1.5\times10^{-3}\,\Omega\cdot\text{cm}=1.5\times10^{-5}\,\Omega\cdot\text{m}$：
\begin{equation}
\boxed{\mu_H=\frac{1.28\times10^{-5}}{1.5\times10^{-5}}\approx0.853\,\text{m}^2/(\text{V}\cdot\text{s})=8530\,\text{cm}^2/(\text{V}\cdot\text{s}).}
\end{equation}
\end{derivationbox}

\begin{formulabox}[答案汇总]
\begin{center}
\begin{tabular}{ll}
\toprule
\textbf{物理量} & \textbf{数值} \\
\midrule
(1) 霍尔系数 $R_H$ & $-1.28\times10^{-5}\,\text{m}^3/\text{C}$ \\
(2) 载流子浓度 $n$ & $4.88\times10^{17}\,\text{cm}^{-3}$ \\
(3) 霍尔迁移率 $\mu_H$ & $0.853\,\text{m}^2/(\text{V}\cdot\text{s})=8530\,\text{cm}^2/(\text{V}\cdot\text{s})$ \\
\bottomrule
\end{tabular}
\end{center}
\end{formulabox}

\begin{insightbox}[物理图像]
\begin{itemize}
    \item 霍尔效应是测量半导体载流子浓度和迁移率的经典方法。霍尔电压的符号直接指示导电类型（n型为负，p型为正）。
    \item 霍尔迁移率 $\mu_H$ 与电导迁移率 $\mu_c$ 通常相差一个霍尔因子 $r_H$（$\mu_H=r_H\mu_c$）。对于简单能带和晶格散射 $r_H\approx1.18$；对于电离杂质散射 $r_H\approx1.93$。
    \item GaAs的电子迁移率极高（可达 $8000\,\text{cm}^2/\text{V}\cdot\text{s}$ 以上），本题计算结果与典型值一致，体现了GaAs作为高迁移率半导体的特性。
\end{itemize}
\end{insightbox}
